\chapter{Markov Chains, Mixing, and Sampling}
\label{ch:stage6-markov-sampling}

This supplement organizes the Markov-chain tools used throughout the book.
The existing random-walk, conductance, coupling, approximate-counting, and
volume-estimation chapters remain in place; this chapter adds the standard
examples and makes the assumptions behind sampling claims explicit.

\section{Mixing and spectral gap}

For a finite irreducible and aperiodic chain with stationary distribution
$\pi$, define
\[
 \tau_x(\varepsilon)=\min\{t:\|P^t(x,\cdot)-\pi\|_{\mathrm{TV}}\leq\varepsilon\},
 \qquad \tau(\varepsilon)=\max_x\tau_x(\varepsilon).
\]
For a reversible chain, the spectral gap is
$\gamma=1-\lambda_*$, where $\lambda_*$ is the largest absolute value of a
nontrivial eigenvalue.  A typical bound has the form
\[
 \tau(\varepsilon)\leq
 \frac{1}{\gamma}\log\frac{1}{\varepsilon\,\pi_{\min}},
 \qquad \pi_{\min}=\min_x\pi(x).
\]
The absolute value and laziness convention matter: a periodic chain can have
an apparent spectral gap but fail to converge in total variation.

\section{Strong stationary times}

\begin{definition}
A stopping time $T$ is a strong stationary time if, for every state $y$ and
every $t$, the event $\{T=t,X_T=y\}$ has probability
$\Pr[T=t]\pi(y)$.
\end{definition}

If $T$ is a strong stationary time, then
\[
 \|\mathcal{L}(X_t)-\pi\|_{\mathrm{TV}}\leq\Pr[T>t].
\]
This is useful when a chain has a natural ``all information has arrived''
time, even when direct eigenvalue calculations are inconvenient.

\section{Coupling and path coupling}

A coupling runs two copies of the same chain jointly.  Once they meet, force
them to use identical future randomness.  The coupling inequality gives
\[
 \|P^t(x,\cdot)-P^t(y,\cdot)\|_{\mathrm{TV}}
 \leq\Pr[X_t\neq Y_t].
\]
Path coupling reduces the analysis to neighboring states when the state space
has a metric: if neighboring copies contract in expected distance by a
factor $\beta<1$ in one step, then arbitrary states contract after extending
the coupling along a shortest path.

\section{Canonical examples}

\subsection{The lazy hypercube walk}

At each step, choose a coordinate uniformly and flip it with probability
$1/2$; otherwise stay put.  The uniform distribution on $\{0,1\}^n$ is
stationary and the chain is reversible.  Its mixing time is
$O(n\log(n/\varepsilon))$; the coupon-collector intuition explains the
$n\log n$ term, since every coordinate must be refreshed.

\subsection{Random-transposition shuffle}

Choose two positions uniformly and transpose them, with a laziness step if
needed.  The uniform distribution on permutations is stationary.  Coupling
or representation-theoretic methods show a mixing scale of order
$n\log n$; this is another example where a local random update eventually
refreshes a global object.

\subsection{Glauber dynamics}

Glauber dynamics updates one local variable at a time according to its
conditional distribution under a target Gibbs measure.  The chain is useful
for colorings, independent sets, and spin systems.  Rapid mixing is not
automatic: it depends on parameters such as temperature, maximum degree, and
the target measure.  A theorem must state the regime in which conductance,
coupling, or a spectral argument applies.

\section{Metropolis--Hastings}

Let $q(x,y)$ be a proposal kernel and let $\pi$ be a target distribution.
From $x$, propose $y\sim q(x,\cdot)$ and accept with probability
\[
 a(x,y)=\min\left\{1,\frac{\pi(y)q(y,x)}{\pi(x)q(x,y)}\right\};
\]
otherwise remain at $x$.

\begin{theorem}[Metropolis--Hastings reversibility]
If the proposal ratios are defined whenever a transition is possible, the
Metropolis--Hastings chain satisfies detailed balance with respect to $\pi$.
Consequently, $\pi$ is stationary.
\end{theorem}

\begin{proof}
For $x\neq y$, the flow in either direction is
\[
 \pi(x)q(x,y)a(x,y)
 =\min\{\pi(x)q(x,y),\pi(y)q(y,x)\},
\]
which is symmetric in $x$ and $y$.  The self-loop probability completes the
transition matrix.
\end{proof}

Stationarity is not rapid mixing.  Mixing still requires irreducibility,
aperiodicity, and a quantitative conductance, coupling, or spectral bound.

\section{Randomized 2-SAT and local chains}

The random-walk algorithm for 2-SAT chooses an unsatisfied clause and flips a
uniformly selected literal.  When a satisfying assignment exists, the Hamming
distance to a fixed satisfying assignment is a supermartingale with an
absorbing target, yielding a polynomial expected hitting-time bound after
appropriate restarts.  This is a hitting-time argument, not a claim that the
walk has the uniform satisfying-assignment distribution.

\section{Graph-coloring chains}

The single-site coloring chain chooses a vertex and recolors it uniformly
from available colors.  Proper colorings are stationary under the uniform
distribution when the chain is irreducible.  Irreducibility can fail when
there are too few colors, and rapid mixing requires a separate condition,
typically a sufficiently large multiple of the maximum degree.  This makes
graph coloring a useful example of the difference between existence,
stationarity, and efficient sampling.

\section{Sampling and counting}

A rapidly mixing chain can produce approximate samples from a target
distribution.  To turn samples into a count, one usually needs a known
reference state, a sequence of overlapping distributions, or a ratio chain.
The variance of the estimator and the burn-in error must both be included.
Conversely, an approximate counter can sometimes support almost-uniform
sampling through self-reduction, but membership testing and conditioning must
be efficient.

\begin{proposition}[Burn-in plus sampling]
If the chain is within $\eta$ of stationarity after $T$ steps and an estimator
based on one stationary sample fails with probability at most $\delta$, then
the same estimator after $T$ burn-in steps has failure probability at most
$\delta+\eta$ for any event bounded between zero and one.
\end{proposition}

This simple accounting prevents a common student error: treating correlated,
not-yet-mixed samples as independent stationary observations.

\section{Student checklist}

For every MCMC or random-walk claim, students should state:

\begin{enumerate}
\item the state space and transition rule;
\item the target stationary distribution;
\item irreducibility and aperiodicity;
\item the mixing metric and error parameter;
\item the proof tool: coupling, conductance, spectral gap, or stopping time;
\item burn-in, sample correlation, estimator variance, and total runtime.
\end{enumerate}

\begin{problem}[Hypercube mixing]
Use a coupon-collector argument to obtain an $O(n\log(n/\varepsilon))$
upper bound for the lazy hypercube walk.
\end{problem}

\begin{problem}[Metropolis--Hastings]
Verify detailed balance for a symmetric proposal kernel and a target Gibbs
distribution.  Identify when the chain becomes reducible.
\end{problem}

\begin{problem}[Path coupling]
Define a distance for a local-update chain and determine a parameter regime in
which neighboring configurations contract in expectation.
\end{problem}

\begin{problem}[Counting from sampling]
Design a ratio estimator using a sequence of nested state spaces.  Bound its
total relative error when each ratio is estimated independently.
\end{problem}
